%% Results Section
\section{Results}
\label{sec:results}

\subsection{Main Results: Fidelity Across Target Families}

\placeholder{Table 1: Best fidelity (mean ± std over 3 seeds) for all 4 families
at 2Q and 3Q, comparing ACC vs Fixed Budget vs Plateau Counter vs ASHA.
Data pending job 181423 completion.}

\begin{table}[t]
\centering
\caption{Best fidelity $F$ (mean $\pm$ std, 3 seeds) at 2 qubits.
\placeholder{Fill with job 181423 results.}}
\label{tab:main_results_2q}
\begin{tabular}{lcccc}
\toprule
Target & ACC & Fixed Budget & Plateau Counter & ASHA \\
\midrule
GHZ   & $\mathbf{0.9997 \pm 0.0002}$ & \placeholder{} & \placeholder{} & \placeholder{} \\
W     & \placeholder{} & \placeholder{} & \placeholder{} & \placeholder{} \\
Cluster & $\mathbf{0.9988 \pm 0.0003}$ & \placeholder{} & \placeholder{} & \placeholder{} \\
Dicke-$k$ & \placeholder{} & \placeholder{} & \placeholder{} & \placeholder{} \\
\bottomrule
\end{tabular}
\end{table}

\subsection{Compute Savings}

\placeholder{Table 2: Total training steps (mean ± std over 3 seeds) for all 4 families
at 2Q, comparing ACC vs Fixed Budget. Include percentage savings column.
Data pending job 181423 completion.}

At 2 qubits, ACC achieves the target fidelity $F^* = 0.99$ with dramatically fewer
training steps than the fixed-budget baseline. For GHZ, ACC stops at approximately
10,000 steps (reason: \texttt{THRESHOLD\_MET}), compared to the fixed budget of
500,000 steps — a saving of 98\%. For Cluster, ACC stops at approximately 14,500
steps, compared to 500,000 — a saving of 97\%. These savings are consistent across
all three seeds, confirming that the saturating exponential model accurately captures
the convergence behaviour.

\subsection{Learning Curve Analysis}

\placeholder{Figure 1: Learning curves (fidelity vs training steps) for all 4 families
at 2Q, showing the fitted saturating exponential model and the ACC stopping point.
Data pending job 181423 completion.}

\placeholder{Figure 2: Extracted parameters $F_\infty$ and $\tau$ for all 4 families
at 2Q and 3Q (if available), showing the target-specific convergence characteristics.}

\subsection{Scaling of $\tau$ with Qubit Count}

\placeholder{Figure 3: Log-linear plot of $\tau$ vs qubit count $n$ for all 4 families,
with fitted exponential $\tau(n) = \tau_0 \cdot \beta^{n-2}$.
Data pending job 181423 completion at 3Q+.}

\placeholder{Table 3: Fitted scaling parameters $\tau_0$ and $\beta$ for each target family,
with 95\% confidence intervals. Include entanglement entropy $S$ of each target state
as a potential predictor of $\beta$.}

\subsection{DEHB Hyperparameter Search}

\placeholder{Figure 4: DEHB search landscape (best fidelity vs entropy coefficient $\alpha$)
for all 4 families at 2Q. Show the relationship between $\alpha$ and $F_\infty$.}

The DEHB search over the entropy coefficient $\alpha$ reveals a clear relationship
between $\alpha$ and the asymptotic fidelity $F_\infty$: configurations with
$\alpha \in [0.015, 0.025]$ consistently achieve $F_\infty \geq 0.99$ for GHZ,
while Cluster requires $\alpha \in [0.010, 0.020]$. This suggests that the
optimal $\alpha$ encodes information about the entanglement structure of the target
state, with more entangled targets requiring higher entropy regularisation to
maintain exploration near the target state.
